\section{Method}
\label{sec:method}

The diagnostic is a single matrix --- the outer product of per-seed
gradients --- and the eigendecomposition that reads its identifiability
content.

\subsection{Setup}
\label{sec:method:setup}

A differentiable agent-based model is a function
$f_\theta : \Xi \to \mathbb{R}^T$ that maps parameters
$\theta \in \mathbb{R}^P$ and a per-seed random key $\xi \in \Xi$ to
an observable trajectory $x = f_\theta(\xi)$. We calibrate against a
reference trajectory tensor $Y_{\mathrm{ref}} \in \mathbb{R}^{M \times T}$,
held fixed throughout, under the unbiased squared maximum-mean-discrepancy
loss with median-heuristic bandwidth \citep{gretton_mmd_2012},
$L(\theta) = \widehat{\mathrm{MMD}}^2\!\bigl(X_\theta, Y_{\mathrm{ref}}\bigr)$,
where $X_\theta = (f_\theta(\xi_m))_{m=1}^{M}$ are $M$ simulator
replicas under independent keys. The per-seed gradient is the
chain-rule contribution from one simulator seed,
\begin{equation}
\label{eq:gm}
g_m \;=\; M \cdot \frac{\partial L}{\partial x_m}\,\frac{\partial x_m}{\partial \theta}
\;\in\; \mathbb{R}^{P},
\qquad
\bar g \;=\; \frac{1}{M}\sum_m g_m \;=\; \nabla_\theta L,
\end{equation}
obtained by a single vector--Jacobian product through the simulator.
The calibrator computes the $g_m$ on its way to $\bar g$; the
diagnostic uses the same $g_m$ a second time, at no marginal cost.
Figure~\ref{fig:method} summarises the pipeline.

\begin{figure}[t]
  \centering
  \begin{tikzpicture}[
    >=Stealth,
    node distance=0.45cm and 0.7cm,
    every node/.style={font=\small},
    box/.style={draw, rounded corners=2pt, fill=gray!8,
                minimum width=2.4cm, minimum height=0.85cm, align=center,
                inner sep=3pt},
    op/.style={font=\scriptsize\itshape, midway, above, sloped,
                yshift=1pt},
    arr/.style={->, line width=0.6pt}
]

\node[box] (theta)   {$\theta \in \mathbb{R}^P$};
\node[box, right=of theta]      (X)    {trajectories \\ $X \in \mathbb{R}^{M \times T}$};
\node[box, right=of X]          (L)    {scalar loss \\ $L = \widehat{\mathrm{MMD}}^2(X, Y_{\mathrm{ref}})$};

\node[box, below=of theta]      (gm)   {per-seed grads \\ $\{g_m\} \in \mathbb{R}^{M \times P}$};
\node[box, right=of gm]         (F)    {OPG matrix \\ $\widehat{F} = \tfrac{1}{M}\sum_m g_m g_m^\top$};
\node[box, right=of F]          (eig)  {eigenpairs \\ $(\lambda_k, v_k)_{k=1}^{P}$};

\draw[arr] (theta) -- (X)  node[op] {vmap simulate};
\draw[arr] (X) -- (L)      node[op] {MMD$^2$};
\draw[arr] (L.south)  -- ++(0,-0.45) -| (gm.north)
                            node[op, pos=0.25] {per-seed vjp};
\draw[arr] (gm) -- (F)     node[op] {outer-product};
\draw[arr] (F) -- (eig)    node[op] {eigh};

\node[font=\scriptsize, below=0.15cm of eig, align=center, text width=2.6cm]
    {stiff $v_1$ \dots\ sloppy $v_P$};
\end{tikzpicture}

  \caption{The OPG diagnostic pipeline. Per-seed gradients computed by vector--Jacobian product through the simulator are the natural inputs for the outer-product matrix $\widehat F$; its eigendecomposition names the parameter combinations the data constrains. The eigendecomposition step is a free byproduct of the gradients the calibrator already needs.}
  \label{fig:method}
\end{figure}

\subsection{The OPG matrix and a Gauss--Newton reading}
\label{sec:method:opg}

The per-seed gradient is the natural unit of identifiability
information already present in the calibrator; the second moment is
the smallest object that exposes its covariance. We define the
\emph{outer-product-of-gradients} (OPG) matrix as
\begin{equation}
\label{eq:opg}
\widehat{F}(\theta) \;=\; \frac{1}{M} \sum_{m=1}^{M} g_m\, g_m^{\top}.
\end{equation}
For an MMD U-statistic loss, the residual structure of the unbiased
estimator \citep{gretton_mmd_2012} licenses interpreting $\widehat F$
as a stochastic generalised Gauss--Newton (GGN) approximation of the
loss Hessian near a minimum, independently of any likelihood
interpretation. \S\ref{sec:intro} sets out the consequence: read this
way, $\widehat F$ is a diagnostic rather than a Fisher surrogate, and
the failure mode \citet{kunstner_opg_2019} identified for
curvature-adaptive optimisers appears here too
(\S\ref{sec:results:predicts-adam}). Throughout we call $\widehat F$
the OPG matrix, never the empirical Fisher.

\subsection{Non-circular validation via non-MMD discrepancies}
\label{sec:method:falsification}

The eigendecomposition $\widehat F\, v_k = \lambda_k\, v_k$ orders the
parameter combinations by how strongly the MMD loss constrains them:
the stiffest direction $v_1$ is the combination the data sees
clearest, the sloppiest direction $v_P$ is the one it cannot
constrain. To rule out the null --- that $\widehat F$ identifies an
MMD-kernel artefact rather than identifiability of the model --- we
perturb $\theta^* \pm \alpha\, v_k$ for fixed $\alpha = 10^{-2}$ and
measure the aggregate discrepancy
\begin{equation}
\label{eq:disc}
\Delta_{\phi}\!\bigl(v_k\bigr) \;=\; \sum_i \bigl| \phi_i\bigl(X_{\theta^* + \alpha v_k}\bigr) - \phi_i\bigl(X_{\theta^*}\bigr) \bigr|
\end{equation}
for three statistics $\phi$ the calibrator never saw: the first four
moments, the autocorrelation up to lag~20, and the four tail quantiles
$\{1\%, 5\%, 95\%, 99\%\}$. If $\widehat F$ identifies real
identifiability, the stiff/sloppy ratio
$\Delta_\phi(v_1) / \Delta_\phi(v_P)$ should be large on every
channel, on every model, with no MMD or kernel embedding appearing
anywhere in the test. Section~\ref{sec:results:falsification} reports
the result on three models.
